\scanpage{169}
\mspage{383}

pour le moment d'ignorer cette possibilité de normalisation -- ce qui met mieux en évidence, en effet, le caractère crucial, dans tous ces calculs $l$-adiques, du rôle joué par les quatre quantités $c,d,t,u\in\widehat{\mathbb Z}$.

\scanpage{170}
\mspage{384}

\uline{45} \uline{L'homomorphisme} $\widetilde M_{0,3}\longrightarrow\mathrm{Gl}(2,\widehat{\mathbb Z})'$, \uline{construction de l'extension} \deleted{\illegible{1}} $\widetilde N_{1,1}$ \uline{de} $\widetilde\Gamma$ \uline{par} $\mathrm{Sl}(2,\mathbb Z)^\wedge$.

On désigne par $\widetilde M_{0,3}$ le groupe des automorphismes \deleted{extérieurs} à lacets de $\unsure{\widehat{\mathcal G}_{0,3}^{+}}$ qui \deleted{\unsure{commutent}} \inserted{induisent} extérieurement \deleted{à l'action de $\unsure{\mathfrak S_3|\mathcal G_3}$, \unsure{qui est}} \inserted{sur le quotient $\unsure{\mathfrak S_3|\mathcal G_3}\simeq\widehat{\mathcal G}_{0,3}^{+}/\widehat\Pi_{0,3}$ l'action triviale -- on a} donc une \inserted{structure d'} extension
\[
(1)\qquad 1\longrightarrow\unsure{\widehat{\mathcal G}_{0,3}^{+}}\longrightarrow\widetilde M_{0,3}\longrightarrow\widetilde N_{0,3}\longrightarrow 1,
\]
\margin{(On désignera par $\widetilde M_{0,3}^{!}$ le sous-groupe formé des \unsure{autom.} qui induisent \unsure{str.\textsuperscript{t}} \unsure{l'ident.\textsuperscript{té}} dans $\unsure{\mathfrak S_3|\mathcal G_3}$ -- il correspond à $\widetilde N_{0,3}^{!}$ \unsure{-- autom.} ext. à lacets qui fixent les classes de lacets (et \unsure{\deleted{induisent}|induisent} l'identité sur $\unsure{\mathfrak S_3|\mathcal G_3}$ \unsure{ext.\textsuperscript{t}}))}

$\widetilde N_{0,3}$ désigne le groupe des \unsure{autom.} \uline{extérieurs} de $\widehat{\mathcal G}_{0,3}^{+}$ qui \deleted{\illegible{3}} \inserted{induisent l'identité dans} $\unsure{\mathfrak S_3|\mathcal G_3}$. On notera que
\[
\Gamma=\Gamma_{\mathbb Q}\hookrightarrow\widetilde\Gamma,
\]
et que l'image inverse de $\Gamma_{\mathbb Q}$ dans l'ext (1) n'est autre que $M_{0,3}$ -- mais pour le moment je n'ai pas de caractérisation bien plausible à proposer de $\Gamma_{\mathbb Q}$ comme sous-groupe de $\widetilde\Gamma$. On a un scindage canonique de l'extension (1), via le sous-groupe $\widetilde\Gamma_{\sigma}$ de $M$ formé des \unsure{autom.} $u\in\widetilde M_{0,3}$ qui commutent à $\sigma_\infty$,



\scanpage{171}

de sorte que $\widetilde M_{0,3}$ se présente comme un produit $\frac12$ direct
\[
(3)\qquad \widetilde M_{0,3}\simeq\underbrace{\widehat{\mathcal G}_{0,3}^{+}}_{\mathrm{Sl}(2,\mathbb Z)'{}^\wedge}\cdot\widetilde\Gamma_\sigma.
\]

Considérons l'homomorphisme canonique
\[
(4)\qquad \unsure{\widehat{\mathcal G}_{0,3}^{+}}\simeq\mathrm{Sl}(2,\mathbb Z)'{}^\wedge\xrightarrow{\varphi_0'}\mathrm{Sl}(2,\widehat{\mathbb Z})'\subset\mathrm{Gl}(2,\widehat{\mathbb Z})'
\]
et l'\deleted{homomorphisme}\inserted{application}
\[
\begin{aligned}
(5)\qquad \widetilde\Gamma_\sigma&\xrightarrow{\varphi_\sigma'}\mathrm{Gl}(2,\widehat{\mathbb Z})'\\
\underline u&\longmapsto\mu(t\sigma_\infty+u)\tau_\infty^\varepsilon=\mu\begin{pmatrix}u&\varepsilon t\\-t&\varepsilon u\end{pmatrix},
\end{aligned}
\]
où \deleted{$\unsure{t,u,\alpha=\alpha_2}$ \illegible{2}}\deleted{\unsure{$\varepsilon=\varepsilon_2$}}
\[
t,u\in\widehat{\mathbb Z},\quad\inserted{\mu\in\widehat{\mathbb Z}^{*},}\quad\varepsilon=\varepsilon_2\in\{\pm1\},\quad\alpha=\alpha_2\in\mathbb Z/2\mathbb Z
\]
sont les quantités (dépendantes de $\underline u\in\widetilde\Gamma_\sigma$) envisagées au \S\ précédent, avec
\[
\varepsilon=(-1)^\alpha.
\]
S'il y a lieu, on écrira $t(\underline u)$, $\underline u(u)$, $\unsure{\alpha(\underline u)}$, $\varepsilon(\underline u)$ pour indiquer la dépendance \unsure{p.r.} à $\underline u$.

\deleted{\uline{Théorème}}\inserted{Considérons l'application (continue)}
\[
(6)\qquad\varphi':\widetilde M_{0,3}\longrightarrow\mathrm{Gl}(2,\widehat{\mathbb Z})'
\]
défini par
\[
(7)\qquad\varphi'(a.\underline u)=\varphi_0'(a)\varphi_\sigma'(\underline u)\qquad\left\{\begin{aligned}&a\in\widehat{\mathcal G}_{0,3}^{+}\\&\underline u\in\widetilde\Gamma_\sigma\end{aligned}\right.
\]

\scanpage{172}
\mspage{385}

\uline{Théorème} L'application $\varphi'$ est un homomorphisme de groupes.

Cela signifie, compte tenu que $\varphi_0$ est déjà un hom de groupes, le théorème signifie que

a) $\varphi_\sigma'$ est un hom. de groupes, i.e.
\[
(8)\qquad\varphi_\sigma'(\underline u'\underline u)=\varphi_\sigma'(\underline u')\unsure{\,\varphi_\sigma'(\underline u)}\qquad(\underline u,\underline u'\in\widetilde\Gamma_\sigma)
\]
b) on a, pour $a\in\unsure{\mathcal G_{0,3}^{+}|\widehat{\mathcal G}_{0,3}^{+}}$, $\underline u\in\widetilde\Gamma_\sigma$
\[
(9)\qquad\varphi_0'\bigl(\underbrace{\underline u a\underline u^{-1}}_{\underline u(a)}\bigr)=\varphi_\sigma'(\underline u)\varphi_0'(a)\varphi_\sigma'\bigl(\underline u^{\deleted{\unsure{-1}}}\bigr)^{-1}.
\]

\deleted{Pour vérifier la}\inserted{Mais la} formule (9) est un fait \inserted{déjà établi,} \deleted{\unsure{conséquence de (8)}, \illegible{1} la} puisque la \deleted{\illegible{1}} construction même de la matrice $\varphi_\sigma'(\underline u)=\mu(t\sigma_\infty+u)\tau_\infty^{\unsure{\varepsilon}}$ aux paragraphes précédents était faite de façon à \deleted{rendre} \unsure{ce} que cette formule soit \unsure{vérifiée} (cf (67)(68), et le cor. au th. (77)\ldots). Pour vérifier (8), \deleted{il faut} \inserted{nous allons} d'abord expliciter les quantités $g,h,f$ pour un produit de deux opérateurs $\unsure{u',u''|u',u}\in\widetilde M_{0,3}$, en termes de leurs valeurs pour $u',u$ eux-mêmes. Un calcul immédiat donne, en ce sens, des ``relations de cocycle''
\[
(10)\qquad\left\{\begin{aligned}
g(\underline u'\underline u)&=\unsure{\underline u'}(g(\underline u)).g(\underline u')\\
h(\underline u'\underline u)&=\underline u'(h(\underline u)).h(\underline u')\\
f(\underline u'\underline u)&=\unsure{f(\underline u')}.\underline u'(f(\underline u))
\end{aligned}\right.
\]

\scanpage{173}

\inserted{\uline{Remarques} 1)} On aura d'ailleurs de $\overline{\mathrm m}$, avec les notations du \S\ précédent
\[
(11)\qquad\left\{\begin{aligned}
\alpha(u'u)&=u'(\alpha(u)).\alpha(u')\\
\beta(u'u)&=u'(\beta(u)).\beta(u')
\end{aligned}\right.
\]

\deleted{\unsure{(NB}} 2) Les formules (10) sont un peu plus \unsure{sympa|sympas} pour $f$, que pour $g$ et $h$, \inserted{$\alpha$ et $\beta$} (-- on aurait envie de changer encore de notation, et de remplacer également $g$ et $h$ par $g^{-1},h^{-1}$ (comme précédemment \unsure{un $g_0$} était devenu $f^{-1}$\ldots).

Désignons encore par les $\overline{\mathrm m}$ lettres $g(u),h(u)$ etc les images de ces quantités dans $\mathrm{Sl}(2,\widehat{\mathbb Z})'$, \deleted{les formules \unsure{(10)} \inserted{(11)} \unsure{donnent}} les formules analogues dans $\mathrm{Sl}(2,\widehat{\mathbb Z})'$ (ou $(\mathrm{Gl}(2,\widehat{\mathbb Z})^{\pm1})'$ pour \deleted{\unsure{(11)}} les $\alpha$ de (11)) qui s'écrivent
\[
(12)\qquad\left\{\begin{aligned}
g(\underline u'\underline u)&=\mathrm{int}(\varphi_\sigma'(\underline u'))(g(\underline u)).g(\underline u')\\
h(\underline u'\underline u)&=\mathrm{int}(\varphi_\sigma'(\underline u'))(h(\underline u)).h(\underline u')\\
f(\underline u'\underline u)&=\deleted{\unsure{\mathrm{int}(\varphi_\sigma'(\underline u'))}}f(\underline u').\bigl(\mathrm{int}\unsure{(\varphi_\sigma(\underline u'))}(f(\underline u))\bigr)\\
\inserted{\alpha(\underline u'\underline u)}&\inserted{{}=\mathrm{int}\,\varphi_\sigma'(\underline u')(\alpha(\underline u)).\alpha(\underline u')}\\
\inserted{\beta(\underline u'\underline u)}&\inserted{{}=\mathrm{int}\,\varphi_\sigma'(\underline u')(\beta(\underline u)).\beta(\underline u')}
\end{aligned}\right.
\]
dans $\mathrm{Gl}(2,\widehat{\mathbb Z})'$.

\deleted{Pour vérifier}\inserted{Ces formules} dans $\mathrm{Gl}(2,\widehat{\mathbb Z})'$ ont bien un caractère entièrement explicite du point de vue calcul effectif, contrairement

\scanpage{174}
\mspage{386}

au cas des formules initiales dans $\widehat{\Pi}_{0,3}$ ou $\unsure{\widehat{\widetilde{\Pi}}_{0,3}}$, où des quantités comme $\underline{u}'(\underbrace{g(\underline{u})}_{g})$ ne sont pas explicitables \deleted{\illegible{1}} \unsure{par} des formules Alg. simples en fonction des données \inserted{$\unsure{g,h,f,}$} $g',h',f'$ \deleted{\unsure{resp.}} $=g(\underline{u})$ etc qui déterminent $\underline{u}$ et $\underline{u}'$.

Ces formules se complètent, pour mémoire, par
\[
\tag{13}
\mu(\underline{u}'\underline{u})=\mu(\underline{u}')\mu(\underline{u})
\]
impliquant notamment
\[
\tag{14}
\varepsilon_2(\underline{u}'\underline{u})=\varepsilon_2(\underline{u}')\varepsilon_2(\underline{u}).
\]

Pour faire le calcul de $\varphi_{\sigma}(\underline{u}'\mathbin{.}\underline{u})$ et vérifier qu'il est égal à $\varphi_{\sigma}(\underline{u}')\varphi_{\sigma}(\underline{u})$, on \deleted{\unsure{aura}} introduira les coefficients $r,s,t,u$ de $f=f(\underline{u})$, $r',s',t',u'$ de \inserted{$f'=$} $f(\underline{u}')$, et $r'',s'',t'',u''$ de $f''=f(\underline{u}'\underline{u})$, de $\overline{\mathrm m}$ \inserted{$\mu,\mu',\mu''$,} $\varepsilon,\varepsilon',\varepsilon''$, d'où
\[
\tag{15}
\deleted{\text{et }\unsure{\text{relations}}}\qquad
\mu''=\mu\mu',\quad \varepsilon''=\varepsilon\varepsilon'
\]
et la formule en $f$ dans (12) s'écrit donc
\[
\tag{16}
f''=f'\underline{u}'f\underline{u}'^{-1}
\]
où on a écrit, comme au \S\ précédent, $\underline{u}'$ au lieu de $\varphi_{\sigma}(\underline{u}')$, \deleted{\unsure{En}} En termes explicites
\[
\begin{pmatrix}r''&s''\\t''&u''\end{pmatrix}
=\begin{pmatrix}r'&s'\\t'&u'\end{pmatrix}
\unsure{\mu'}\begin{pmatrix}u'&\varepsilon't'\\-t'&\varepsilon'u'\end{pmatrix}
\begin{pmatrix}r&s\\t&u\end{pmatrix}
\unsure{\deleted{\frac{\mu'}{\mu'}}|\frac{\mu'}{\mu'}}
\begin{pmatrix}\varepsilon'u'&-\varepsilon't'\\t'&u'\end{pmatrix}
\]
\margin{NB $\det\underline{u}'=\mu'$}

\scanpage{175}

Le produit des trois derniers termes du second membre n'est pas engageant :
\[
\underline{u}'f\underline{u}'^{-1}
=\mu'\deleted{\unsure{\frac{1}{\mu'}}}
\begin{pmatrix}
\varepsilon'ru'^2+st'u'+tt'u'+\varepsilon'ut'^2
&-\varepsilon'rt'u'+su'^2-tt'^2+\varepsilon'uu't'\\
-\varepsilon'ru't'\unsure{-}st'^2+tu'^2+\varepsilon'uu't'
&-\varepsilon'rt'^2\unsure{-}su't'-tu't'+\varepsilon'uu'^2
\end{pmatrix},
\]
mais le produit par $f'$ se simplifie miraculeusement, grâce à la relation
\[
\unsure{t'^2}+u'^2=\varepsilon'/\mu'
\]
et on trouve
\[
\deleted{\unsure{(17)}}\qquad
f''=\begin{pmatrix}r''&s''\\t''&u''\end{pmatrix}
=\deleted{\unsure{\frac{1}{\mu'}}}
\begin{pmatrix}
\unsure{\varepsilon|*}\ \text{bordel}&\text{bordel}\\
\varepsilon'tu'+ut'&-\varepsilon'tt'+uu'
\end{pmatrix}
\]
d'où
\[
\tag{18}
t''=\frac{1}{\mu'}(\varepsilon'tu'+ut')
\qquad
u''=\frac{1}{\mu'}(-tt'+uu')
\]
et par suite (tenant compte de \inserted{$\mu''=\mu\mu'$,} $\varepsilon''=\varepsilon\varepsilon'$)
\[
\deleted{\unsure{(19)}}\qquad
\underline{u}''=\mu\mu'\deleted{\unsure{\mu'}}
\begin{pmatrix}
-\varepsilon'tt'+uu'&\varepsilon tu'+\varepsilon\varepsilon'ut'\\
-\varepsilon'tu'-ut'&-\varepsilon\varepsilon'tt'+\varepsilon\varepsilon'uu'
\end{pmatrix}
\]
et la formule à vérifier \unsure{$u$},
\[
\underline{u}''=\underline{u}'\underline{u}
\]
se réduit à
\[
\mu\mu'\begin{pmatrix}
-tt'+uu'&\varepsilon tu'+\varepsilon\varepsilon'ut'\\
-\varepsilon'tu'\deleted{\unsure{\varepsilon'}}-ut'&-\varepsilon\deleted{\unsure{\varepsilon'}}tt'+\varepsilon\varepsilon'uu'
\end{pmatrix}
=\mu'\begin{pmatrix}u'&\varepsilon't'\\-t'&\varepsilon'u'\end{pmatrix}
\mu\begin{pmatrix}u&\varepsilon t\\-t&\varepsilon u\end{pmatrix}
\]
qui est bien vraie, OK !.

Ici encore, on sent que ces calculs devraient s'insérer dans un tapis d'algèbre linéaire que j'ai la flemme de dérouler in extenso. Il faudrait paraphraser le groupe \inserted{\unsure{Ou plutôt}} $\widetilde{M}_{0,3}$ -- décrit par ces systèmes $(g,h,f)$ -- en termes \inserted{$(\alpha,\beta,f)$} des réductions de $g,h,f$

\scanpage{176}
\mspage{387}

dans $\mathrm{Gl}(2,\widehat{\mathbb Z})^{\pm1\,\prime}_{\unsure{\mathfrak g}}$, puis en remplaçant $\widehat{\mathbb Z}$ par un anneau quelconque, et bien sûr $\mathrm{Gl}(2,\widehat{\mathbb Z})'$ par $\mathrm{Gl}(2,\unsure{A|\widehat{\mathbb Z}})$ lui-$\overline{\mathrm m}$. Peut-être y reviendrai-je finalement....

Ayant ainsi construit
\[
\varphi:\widetilde{M}_{0,3}\longrightarrow\mathrm{Gl}(2,\widehat{\mathbb Z})'=\mathrm{Gl}(2,\widehat{\mathbb Z})/\mathord{\pm1}
\]
on en déduit, par image inverse de l'extension \inserted{centrale} $\mathrm{Gl}(2,\widehat{\mathbb Z})$ de $\mathrm{Gl}(2,\widehat{\mathbb Z})'$ par \deleted{$\unsure{\mathbb Z}$} $\{\pm1\}$, une extension centrale de $\widetilde{M}_{0,3}$ par $\{\pm1\}$, que je note $\unsure{\widetilde{N}_{1,1}}$, et on a un hom d'extensions
\[
\tag{19}
\diagram{\begin{tikzcd}[ampersand replacement=\&]
1\arrow[r]\&\{\pm1\}\arrow[r]\arrow[d,equal]\&\widetilde{N}_{1,1}\arrow[r]\arrow[d,"\varphi"]\&\deleted{\unsure{\mathcal T}}\widetilde{M}_{0,3}\arrow[r]\arrow[d,"\varphi'"]\&1\\
1\arrow[r]\&\{\pm1\}\arrow[r]\&\mathrm{Gl}(2,\widehat{\mathbb Z})\arrow[r]\&\mathrm{Gl}(2,\widehat{\mathbb Z})'\arrow[r]\&1
\end{tikzcd}}
\]

\deleted{Par ailleurs, désignant par \deleted{$\unsure{\widetilde{N}_{1,1}}$} $\unsure{\widetilde{\mathcal T}_{1,1}^{+}|\widehat{\mathcal T}_{1,1}^{+}}$ l'image inverse de \deleted{$\unsure{\widetilde{N}_{1,1}}$} $\widehat{\mathcal G}_{0,3}^{+}\simeq\unsure{\widehat{\mathrm{Sl}(2,\mathbb Z)}'}\subset\widetilde{M}_{0,3}$ dans l'extension \deleted{$\unsure{\mathcal T}$} $\widetilde{N}_{1,1}$ de $\widetilde{M}_{0,3}$ par $\pm1$, on aura
\[
\diagram{\begin{tikzcd}[ampersand replacement=\&]
1\arrow[r]\&\unsure{\widehat{\mathcal T}_{1,1}^{+}|\widetilde{\mathcal T}_{1,1}^{+}}\arrow[r]\&\widetilde{N}_{1,1}
\end{tikzcd}}
\]
}

\scanpage{177}

Il faudrait vérifier que l'hom $\varphi'$ qu'on vient de construire coïncide sur $M_{0,3}\subset\widetilde{M}_{0,3}$ avec l'hom. ``modulaire'' bien connu du \S\ 39 \unsure{(1)} -- il suffit de le vérifier sur le sous groupe $\Gamma_{\mathbb Q,\sigma}$ de $\widetilde{\Gamma}_{\sigma}$ correspondant à $\Gamma_{\mathbb Q}$ -- ça ne doit pas être bien difficile, et je vais \unsure{l'admettre} ici. Ceci \deleted{\unsure{étant}} admis, on trouve que l'image inverse de $M_{0,3}$ dans \deleted{$\unsure{\mathcal T}$} $\widetilde{N}_{1,1}$ n'est autre (à isom. can. près) que $N_{1,1}$ lui $\overline{\mathrm m}$ -- en tous cas l'image inverse de $\widehat{\mathcal G}_{0,3}^{+}$ \deleted{\unsure{sur}} n'est autre que $\widehat{\mathcal T}_{1,1}^{+}$ \inserted{$\simeq\mathrm{Sl}(2,\mathbb Z)^\wedge$}. On aura alors (avec cette identification)
\[
\tag{20}
\diagram{\begin{tikzcd}[ampersand replacement=\&]
1\arrow[r]\&\begin{gathered}\widehat{\mathcal T}_{1,1}^{+}\\\unsure{\Vert}\ \mathrm{Sl}(2,\mathbb Z)^\wedge\end{gathered}\arrow[r]\arrow[d]\&\deleted{\unsure{\mathcal G}}\widetilde{N}_{1,1}\arrow[r]\arrow[d,"\varphi"]\&\widetilde{\Gamma}\arrow[r]\arrow[d,"\unsure{\chi}"]\&1\\
1\arrow[r]\&\mathrm{Sl}(2,\widehat{\mathbb Z})\arrow[r]\&\mathrm{Gl}(2,\widehat{\mathbb Z})\arrow[r,"\det"]\&\widehat{\mathbb Z}^{*}\arrow[r]\&1
\end{tikzcd}}
\]

\scanpage{178}
\mspage{388}

donc on a \unsure{réconstruit|reconstruit} en termes d'\unsure{automs} à lacets de $\widehat{\Pi}_{0,3}$ le diagramme (1) du \S\ 39, ou plutôt un diagramme \unsure{un} plus grand, dans lequel se plonge celui du \S\ 39 -- il resterait encore à caractériser $\Gamma_{\mathbb Q}$ comme sous-groupe de $\widetilde{\Gamma}$.

Le diagramme déduit directement du th. précédent
\[
\tag{21}
\diagram{\begin{tikzcd}[ampersand replacement=\&]
1\arrow[r]\&\widehat{\mathcal G}_{0,3}^{+}\arrow[r]\arrow[d]\&\widetilde{M}_{0,3}\arrow[r]\arrow[d,"\varphi'"]\&\widetilde{\Gamma}\arrow[r]\arrow[d,"\chi"]\&1\\
1\arrow[r]\&\mathrm{Sl}(2,\widehat{\mathbb Z})'\arrow[r]\&\unsure{\mathrm{Gl}(2,\widehat{\mathbb Z})'}\arrow[r,"\det"']\&\widehat{\mathbb Z}^{\unsure{*}}\arrow[r]\&1
\end{tikzcd}}
\]
peut être considéré comme étant déduit du précédent par réduction mod les \deleted{centres} \inserted{sous groupes centraux} $\{\pm1\}$ dans les deux lignes.

Que devient dans (20) le scindage (3) de $\widetilde{M}_{0,3}$ ? Y-a-t-il une façon privilégiée, pour \unsure{tt} élément de $\widetilde{N}_{1,1}$ -- i.e. un couple d'un $u=(g,h,f)\in\widetilde{M}_{0,3}$, et d'un $\underline{u}_0\in\mathrm{Gl}(2,\widehat{\mathbb Z})$

\scanpage{179}

qui relève l'élément $\underline{u}=\varphi'_{\unsure{\sigma}}(u)$ de $\mathrm{Gl}(2,\widehat{\mathbb Z})'$, de le ``normaliser'' mod \deleted{\illegible{1}} composition avec un élément de $\mathrm{Sl}(2,\mathbb Z)^\wedge$, de façon à ce qu'il ait \deleted{une} (disons) une propriété de commutation avec $\sigma_\infty$, du type $u(\sigma_\infty)=\deleted{\varepsilon}\sigma_\infty$, \deleted{\illegible{1}} $\unsure{u_0|u}(\sigma_\infty)=\pm\sigma_\infty$ ?

Une autre façon de poser la question, en termes du \unsure{s-groupe} $\widetilde{\Gamma}_\sigma\subset\widetilde{M}_{0,3}$, est la suivante : y-a-t-il un relèvement canonique $\widetilde{\Gamma}_\sigma\longrightarrow\widetilde{N}_{1,1}$, associant donc à tout $\underline{u}=(g,f)\in\widetilde{\Gamma}_\sigma$ (\unsure{Càd} un $\underline{u}=\mu(t\sigma_\infty+u)\tau_\infty^\alpha\in\mathrm{Gl}(2,\widehat{\mathbb Z})'$\unsure{,|\}} \deleted{\unsure{mod}} i.e. un $\underline{u}$ défini mod signe) un représentant $\underline{u}_0\in\mathrm{Gl}(2,\widehat{\mathbb Z})$ -- i.e. une façon de lever cette ambiguïté du signe ? En fait, si on remplace $\widetilde{\Gamma}_\sigma$ par $\Gamma_{\mathbb Q,\sigma}$, les \deleted{\illegible{1}}\inserted{\unsure{yogas}} Standard \unsure{suggèrent} que les relèvements en

\scanpage{180}
\mspage{389}

question \unsure{devraient|doivent} correspondre aux courbes elliptiques sur $\mathbb Q$, munies d'une rigidification de Jacobi d'échelon $2$, \unsure{correspondant} (dans l'isomorphisme $\deleted{\unsure{\mathcal T}}\unsure{\mathcal M}_{1,1\,\mathbb Q}'\simeq U_{0,3\,\mathbb Q}$) au point $1/2$ de $U_{0,3}(\mathbb Q)$. L'indétermination pour le choix d'une telle courbe elliptique est justement dans $H^1(\mathbb Q,\unsure{\{\pm1\}})$ (c'est une \unsure{donn.} d'ext. quadratique \unsure{de|à} $\mathbb Q$) qui s'identifie à $\mathrm{Hom}(\Gamma_{\mathbb Q},\deleted{\unsure{\mu}}\{\pm1\})$ -- c'est donc bien aussi l'indétermination pour un relèvement. Ceci prouve \unsure{essentiellement} que l'idée suggérée par le yoga général est bel et bien \unsure{correct|correcte} -- et cela montre en $\overline{\mathrm m}$ temps qu'il n'y a pas ($\overline{\mathrm m}$ sur $\Gamma_{\mathbb Q,\sigma}$) de choix \unsure{vraiment} privilégié. Qu'il y ait de tels relèvements \inserted{(sur $\Gamma_{\mathbb Q,\sigma}$)} provient du fait qu'il y a bel et bien des courbes

\scanpage{181}

elliptiques répondant à la question, la plus simple étant donnée par l'équation
\[
(22)\qquad y^{\unsure{2}}-x(x-1)(x-1/2)=0
\]
-- \unsure{et} voilà le choix fait ! Il resterait à essayer de l'expliciter \deleted{par} dans un \unsure{formulaire} calculatoire, si ce n'est pas \unsure{sans espoir}. Il faudrait que j'examine cette \deleted{question} question de plus près, à l'occasion.

\scanpage{182}
\mspage{390}

\uline{46}\quad \uline{Notations définitives} (?) -- \uline{nouvelle récapitulation}

(\uline{pour} les autom. à lacets de $\widehat{\Pi}_{0,3}$ et de $\widehat{\mathcal G}_{0,3}^{+}\simeq\unsure{\mathrm{Sl}(2,\mathbb Z)^{\prime\wedge}}$).

\textcircled{I}\quad \uline{Les groupes} $\widetilde M_{0,3}$ \uline{et} $\widetilde N_{0,3}^{!}$.

On pose

$(1)\qquad\widetilde M_{0,3}=$ groupe des autom de $\widehat{\mathcal G}_{0,3}^{+}\simeq\unsure{\mathrm{Sl}(2,\mathbb Z)^{\wedge\prime}}$ qui invarient $\widehat{\Pi}_{0,3}$, et qui respectent la ``structure à lacets'' i.e. \unsure{satisfont} $u(\unsure{\varepsilon_i|\varepsilon_0})=\operatorname{int}(f^{-1})(\varepsilon_i^\mu)$, avec $f\in\widehat{\mathcal G}_{0,3}^{+}$, $i\in\{0,1,\infty\}$, $\mu\in\widehat{\mathbb Z}^{*}$ (OBS $f\in\widehat{\Pi}_{0,3}$, et $\unsure{f|f_0}$ est déterminé alors dans $L_0\backslash\widehat{\Pi}_{0,3}$)

$\overset{\mathrm{restr.}}{\overset{\sim}{\longrightarrow}}$ groupe des autom \deleted{ext.} à lacets de $\widehat{\Pi}_{0,3}$ qui commutent extérieurement à l'action extérieure naturelle de $\unsure{\mathfrak S_3|\mathcal G_3}$.

\margin{\uline{NB} $\widetilde M_{0,3}$ est aussi le normalisateur de $\widehat{\mathcal G}_{0,3}^{+}$ dans $\widehat{\widehat{\mathcal G}}_{0,3}$, et de $\overline{\mathrm m}$ $\widetilde N_{0,3}$ le \unsure{normalisateur|normalis.} de $\unsure{\widehat{\mathcal T}_{0,3}}=\unsure{\mathcal T_{0,3}}$ dans $\widehat{\widehat{\mathcal T}}_{0,3}$}

On a une structure d'extension
\[
(2)\qquad 1\longrightarrow\widehat{\Pi}_{0,3}\longrightarrow\widetilde M_{0,3}\longrightarrow\widetilde N_{0,3}\longrightarrow1,
\]
\unsure{où} \deleted{$\unsure{N_0}$}
\[
(3)\qquad\widetilde N_{0,3}\overset{\mathrm{dfn}}{=}\widetilde M_{0,3}/\widehat{\Pi}_{0,3}
\]
$\simeq$ groupe des autom \uline{ext} à lacets de $\widehat{\Pi}_{0,3}$, commutant à $\unsure{\mathfrak S_3|\mathcal G_3}$.

L'hom. naturel : $\unsure{\mathfrak S_3|\mathcal G_3}\longrightarrow\operatorname{Aut}(\unsure{\mathfrak S_3|\mathcal G_3})$ est un \unsure{iso}
\[
(4)\qquad\unsure{\mathfrak S_3|\mathcal G_3}\overset{x\mapsto\operatorname{int}(x)}{\overset{\sim}{\longrightarrow}}\operatorname{Aut}(\unsure{\mathfrak S_3|\mathcal G_3}),
\]
et l'hom. naturel
\[
\widetilde M_{0,3}\longrightarrow\operatorname{Aut}(\unsure{\mathfrak S_3|\mathcal G_3}=\widehat{\mathcal G}_{0,3}^{+}/\widehat{\Pi}_{0,3})
\]
définit alors un homomorphisme
\[
(5)\qquad\widetilde M_{0,3}\longrightarrow\unsure{\mathfrak S_3|\mathcal G_3}
\]
qui peut aussi s'interpréter comme le composé
\[
(6)\qquad\widetilde M_{0,3}\hookrightarrow
\underset{\substack{\text{tous les autom.}\\\text{à lacets de }\widehat{\Pi}_{0,3}}}{\widehat{\widehat{\mathcal G}}_{0,3}}
\xrightarrow{\substack{\text{action sur les}\\\text{classes de lacets}}}
\unsure{\mathfrak S_3|\mathcal G_3},
\]

\scanpage{183}

et qui est trivial sur $\widehat{\Pi}_{0,3}$, donc se factorise en
\[
(7)\qquad\widetilde N_{0,3}\longrightarrow\unsure{\mathfrak S_3|\mathcal G_3}.
\]
Les noyaux de (5), (7), sont notés $\widetilde M_{0,3}^{!}$, $\widetilde N_{0,3}^{!}$ respectivement, de sorte qu'on a un diagramme de suites exactes

\unsure{\diagram{
\begin{tikzcd}[ampersand replacement=\&]
 \& \& \inserted{1}\arrow[d] \& \& \& \\
 \& \& \inserted{\Pi_{0,3}}\arrow[d] \& \& \& \\
(8) \& 1\arrow[r] \& \widetilde M_{0,3}^{!}\arrow[r]\arrow[d] \& \widetilde M_{0,3}\arrow[r]\arrow[d] \& \unsure{\mathfrak S_3|\mathcal G_3}\arrow[r]\arrow[d,equal] \& 1 \\
(9) \& 1\arrow[r] \& \widetilde N_{0,3}^{!}\arrow[r]\arrow[d] \& \widetilde N_{0,3}\arrow[r] \& \unsure{\mathfrak S_3|\mathcal G_3}\arrow[r] \& 1, \\
 \& \& \inserted{1} \& \& \&
\end{tikzcd}
}}

(9) pouvant être considéré comme quotient de (8) par $\widehat{\Pi}_{0,3}$.

\inserted{\unsure{Notons} aussi la structure d'ext. \unsure{correspondante}
\[
(9\mathrm{bis})\qquad1\longrightarrow\widehat{\Pi}_{0,3}\longrightarrow\widetilde M_{0,3}^{!}\longrightarrow\widetilde N_{0,3}^{!}\longrightarrow1
\]}

On a aussi un \unsure{hom.} canonique
\[
(10)\qquad\unsure{\mathfrak S_3|\mathcal G_3}\simeq\widehat{\mathcal G}_{0,3}^{+}/\widehat{\Pi}_{0,3}\hookrightarrow\widetilde N_{0,3},
\]
dont le composé avec 7 est l'identité, d'où un iso. can.
\[
(11)\qquad\widetilde N_{0,3}\simeq\unsure{\mathfrak S_3|\mathcal G_3}\times\widetilde N_{0,3}^{!}
\]
D'autre part, on a aussi une inclusion
\[
(12)\qquad
\underset{\substack{\unsure{\wr|\simeq}\\\unsure{\mathfrak S_3|\mathcal G_3}\times\mathbb Z/2\mathbb Z}}{\mathcal T_{0,3}}
\hookrightarrow\widetilde N_{0,3}
\]
qui sur le facteur $\unsure{\mathfrak S_3|\mathcal G_3}\inserted{=\mathcal T_{0,3}^{+}}$ \unsure{coïncide} avec (10), et qui \deleted{sur} applique l'élément non-trivial $\tau$ de $\mathbb Z/2\mathbb Z\inserted{=\mathcal T_{0,3}^{!}}$ sur l'élément, également noté $\tau$, du deuxième facteur $\widetilde N_{0,3}^{!}$ de (11).

\scanpage{184}
\mspage{391}

L'homomorphisme composé
\[
(13)\qquad \widetilde M_{0,3}\unsure{\hookrightarrow|\subset}
\unsure{\widehat{\widehat{\mathcal G}}_{0,3}}
\xrightarrow[\chi]{\text{\unsure{multiplicateur}}}
\widehat{\mathbb Z}^{*}
\]
\deleted{\unsure{donne|définit}} un hom
\[
(14)\qquad \widetilde M_{0,3}\xrightarrow{\chi}\widehat{\mathbb Z}^{*}
\]
surjectif, se factorisant par $\widetilde N_{0,3}$ et $\overline{\mathrm m}$ \unsure{par} $\widetilde N_{0,3}/\unsure{\mathfrak S_3|\mathcal G_3}$

\diagram{\[
(15)\qquad
\begin{tikzcd}[ampersand replacement=\&]
\widetilde N_{0,3}\arrow[rr,"\chi"]\arrow[dr] \&\& \widehat{\mathbb Z}^{*} \\
\& \begin{gathered}\widetilde N_{0,3}/\unsure{\mathfrak S_3|\mathcal G_3}\\ \simeq\widetilde N_{0,3}^{!}\end{gathered}\arrow[ur] \&
\end{tikzcd}
,\]}

et donnant lieu à des structures d'extensions

\diagram{\[
(16)\qquad
\begin{tikzcd}[ampersand replacement=\&]
1\arrow[r] \& \widetilde M_{0,3}^{+}\arrow[r]\arrow[d] \& \widetilde M_{0,3}\arrow[r,"\chi"]\arrow[d] \& \widehat{\mathbb Z}^{*}\arrow[r]\arrow[d,equal] \& 1 \\
1\arrow[r] \& \widetilde N_{0,3}^{+}\arrow[r] \& \widetilde N_{0,3}\arrow[r,"\chi"] \& \widehat{\mathbb Z}^{*}\arrow[r] \& 1
\end{tikzcd}
,\]}

où on a posé
\[
(17)\qquad
\left\{\begin{aligned}
\widetilde M_{0,3}^{+}&=\operatorname{Ker}(\widetilde M_{0,3}\to\widehat{\mathbb Z}^{*})\\
\widetilde N_{0,3}^{+}&=\operatorname{Ker}(\widetilde N_{0,3}\to\widehat{\mathbb Z}^{*})
\end{aligned}\right.
\]
-- la deuxième ligne de (16) peut s'interpréter comme le quotient de la première ligne par $\widehat\Pi_{0,3}$.

On a un plongement canonique
\[
(18)\qquad
\Gamma_{\mathbb Q}\overset{\mathrm{dfn}}{=}
\operatorname{Gal}\left(
\underbrace{\overline{\mathbb Q}_0}_{\inserted{\substack{\text{clôture}\\\text{alg. de }\mathbb Q\\\text{dans }\mathbb C}}}
/\mathbb Q_0\right)
\unsure{\hookrightarrow|\subset}\widetilde N_{0,3}^{!}
\]
appliquant l'\unsure{autom.} de conjugaison

\scanpage{185}

complexe $\tau$ sur l'élément déjà noté $(\tau)$ (cf (12)). L'image inverse de \unsure{ce|le} sous-groupe $\Gamma_{\mathbb Q}$ \deleted{dans $\unsure{M}$}\inserted{de $\widetilde N_{0,3}^{!}$} par $\widetilde M_{0,3}\to\widetilde N_{0,3}^{!}\simeq\widetilde N_{0,3}/\unsure{\mathfrak S_3|\mathcal G_3}$, i.e l'image inverse de \deleted{$\unsure{\mathfrak S_3|\mathcal G_3}\times\Gamma_{\mathbb Q}\subset$}
\[
(19)\qquad
N_{0,3}\overset{\mathrm{dfn}}{=}
\unsure{\mathfrak S_3|\mathcal G_3}\times\Gamma_{\mathbb Q}
\subset\widetilde N_{0,3}
=\unsure{\mathfrak S_3|\mathcal G_3}\times\widetilde N_{0,3}^{!},
\]
est l'extension \deleted{\unsure{de}} notée $M_{0,3}$ par ailleurs :

\margin{[On a une \deleted{\illegible{1}}\inserted{sous-}extension de \inserted{\unsure{l'ext.}} $\widetilde M_{0,3}$ de (2)]}

\diagram{\[
(20)\qquad
\begin{tikzcd}[ampersand replacement=\&]
1\arrow[r] \& \widehat\Pi_{0,3}\arrow[r]\arrow[d,"\unsure{\iota}"] \& M_{0,3}\arrow[r]\arrow[d,shift right=1.5pt]\arrow[d,no head,shift left=1.5pt,shorten <=7pt,shorten >=7pt] \& \overbrace{\unsure{\mathfrak S_3|\mathcal G_3}\times\Gamma_{\mathbb Q}}^{N_{0,3}}\arrow[r]\arrow[d,"{(\alpha,\beta)\to\beta}"] \& 1 \\
1\arrow[r] \& \widehat{\mathcal G}_{0,3}^{+}\arrow[r] \& M_{0,3}\arrow[r] \& \Gamma_{\mathbb Q}\arrow[r] \& 1
\end{tikzcd}
.\]}

Bien entendu, c'est par $M_{0,3}$ \inserted{et $N_{0,3}$} plutôt que par $\widetilde M_{0,3}$ \inserted{et $\widetilde N_{0,3}$} que nous sommes intéressés -- donc par le sous-groupe remarquable $\Gamma_{\mathbb Q}$ de $\widetilde N_{0,3}^{!}$. Nous allons "déterminer" par le calcul le groupe $\widetilde N_{0,3}^{!}$ -- et il s'agira dans la suite du travail de trouver un maximum de conditions nécessaires sur des éléments de $\widetilde N_{0,3}^{!}$ pour provenir de $\Gamma_{\mathbb Q}$ -- \unsure{jusqu'à} trouver, si possible, une caractérisation (au moins plausible, sinon prouvée !) des éléments de $\Gamma_{\mathbb Q}$ -- i.e. de $\Gamma_{\mathbb Q}$ lui-$\overline{\mathrm m}$.

\scanpage{186}
\mspage{392}

\Needspace{8\baselineskip}
II \uline{Détermination calculatoire des éléments de} $\widetilde N_{0,3}^{!}$.\par\nobreak
1) \uline{Détermination par} $h,g,f$.

\deleted{La \unsure{donnée} de $u\in\unsure{\widetilde M_{0,3}^{!}|\widetilde N_{0,3}^{!}}$}

Rappelons que
\[
\begin{aligned}
(21)\qquad
&\unsure{\widetilde N_{0,3}^{!}|\widetilde M_{0,3}^{!}}
=\text{autom. de }\widehat{\mathcal G}_{0,3}^{+}
\simeq\unsure{\mathrm{Sl}(2,\mathbb Z)^{\wedge\prime}|\widehat{\mathrm{Sl}(2,\mathbb Z)^{\prime}}}
\text{ qui stabilisent}\\
&\widehat\Pi_{0,3},\quad\text{induisent l'identité sur }
\unsure{\mathfrak S_3|\mathcal G_3}
=\widehat{\mathcal G}_{0,3}^{+}/\widehat\Pi_{0,3},\\
&\text{et respectent la structure à lacets (définie}\\
&\text{par la classe de }\varepsilon_0\text{)}
\end{aligned}
\]

Un $u\in\widetilde M_{0,3}^{!}$ est déterminé par un système
\[
(22)\qquad (h,g,f)\quad
\deleted{\unsure{\in(\widehat{\Pi}_{0,3})^{3}}\,\unsure{\times}\,\unsure{\widehat{\Pi}_{0,3}}}
\qquad(h,g,f\in\widehat\Pi_{0,3})
\]
\margin{$f$ déterminé mod $f\mapsto\ell_0^{n}f$, $n\in\widehat{\mathbb Z}$}
\deleted{\unsure{satisf.}}\inserted{via} les formules
\[
(23)\qquad
\left\{\begin{aligned}
u(\sigma_\infty)&=h\sigma_\infty\\
u(\rho)&=g\rho
\end{aligned}\right.
\]
\deleted{$\unsure{u(\varepsilon_0)}=\unsure{\operatorname{int}}\unsure{(}$}

la condition que ces formules définissent un automorphisme (ou du moins un endomorphisme) de $\widehat{\mathcal G}_{0,3}^{+}$ (qui induira néc. l'identité sur $\unsure{\mathfrak S_3|\mathcal G_3}$) est exprimée par les formules
\[
(24)\qquad (h\sigma_\infty)^{2}=(g\rho)^{3}=1,
\]
qui s'écrivent aussi
\[
(25)\qquad
\boxed{\left\{\begin{aligned}
h\sigma_\infty(h)&=1\\
g\rho(g)\rho^{\unsure{2}}(g)&=1
\end{aligned}\right.},
\]
et la condition que $u$ respecte la structure à lacets est exprimée par
\[
(26)\qquad u(\varepsilon_0)=\operatorname{int}(f^{-1})(\varepsilon_0^{\mu})
\qquad \mu\in\widehat{\mathbb Z}^{*}\text{ convenable},
\]
\margin{relation des lacets dans $\widehat{\mathcal G}_{0,3}^{+}$ (forme brute)}
\unsure{équivalente|équivalent} à la condition

\scanpage{187}

\margin{relation des lacets dans $\widehat\Pi_{0,3}$ (forme brute)}
\[
(27)\qquad
u(\ell_0)=\operatorname{int}(f^{-1})(\ell_0^{\mu}),
\]
elle s'explicite aisément en termes de $h,g,f$ (sous forme (26) de préférence) compte tenu que $\varepsilon_0=\sigma_\infty\rho$
\[
\underbrace{(h\sigma_\infty)(g\rho)}_{h\sigma_\infty(g)\underbrace{\sigma_\infty\rho}_{\varepsilon_0}}
=\underbrace{\operatorname{int}(f^{-1})\varepsilon_0^{\mu}}_{f^{-1}\varepsilon_0^{\mu}f},\quad\text{i.e}
\]
\[
f h\sigma_\infty(g)\underbrace{\varepsilon_0 f^{-1}}_{\unsure{\varepsilon_0(f^{-1})}\varepsilon_0=\sigma_\infty(\rho(f)^{-1})\varepsilon_0}
=\varepsilon_0^{\mu}\qquad\text{soit}
\]
\[
f h\sigma_\infty(g\rho(f)^{-1})=\varepsilon_0^{\mu-1}
\qquad\text{soit, posant}
\]
\[
(28)\qquad \unsure{\nu}\overset{\mathrm{dfn}}{=}\frac{\mu-1}{2}\in\widehat{\mathbb Z}
\]
\margin{possible, car $\mu\in\widehat{\mathbb Z}^{*}$ donc $\mu\equiv1\ (2)$}
la relation (compte tenu que $\varepsilon_0^{2}=\ell_0$)
\[
f h\sigma_\infty(g\rho(f)^{-1})=\ell_0^{\nu},
\]
ou encore, appliquant $\sigma_\infty$, comme $\sigma_\infty(\ell_0)=\ell_1$
\[
\sigma_\infty(f)\sigma_\infty(h)g\rho(f)^{-1}=\ell_1^{\nu}
\]
soit enfin
\[
(29)\qquad
\boxed{\sigma_\infty(h)g=\sigma_\infty(f)^{-1}\ell_1^{\nu}\rho(f)}
\]
(\uline{relation des lacets en} $h,g,f$).

\margin{\uline{NB} $\nu$\inserted{$\in\widehat{\mathbb Z}$} est uniquement déterminé par $g,h$, et \unsure{on a} encore $2\nu+1\in\widehat{\mathbb Z}^{*}$ $2\nu+1$ est le multiplicateur\ldots}

\uline{NB} Il faut en plus exiger, pour que $h,g,f$ définissent bien un \uline{autom}, que (24) soit une \uline{présentation-génération} de $\widehat{\mathcal G}_{0,3}^{+}$ -- par la suite, nous sous-entendrons tjs tacitement cette "condition liminaire".
